\documentclass[11pt]{letter}

\usepackage{geometry}
\geometry{margin=1in}
\usepackage{setspace}
\usepackage{hyperref}

\signature{Decory Edwards}
\address{Decory Edwards \\ Johns Hopkins University \\ Department of Economics}

\begin{document}

\begin{letter}{Hiring Committee \\ Demand Side Analytics}

\opening{Dear Hiring Committee,}

I am writing to apply for the Quantitative Analyst position in the Energy Sector at Demand Side Analytics. I am currently a PhD candidate in Economics at Johns Hopkins University. My training combines quantitative analysis, economic reasoning, and data driven empirical methods. I am motivated by problems that require rigorous analysis and clear communication to inform decisions at the intersection of energy, policy, and economics. Demand Side Analytics’ focus on applied analysis of the electric grid and natural gas systems aligns closely with how I approach research.

My research agenda is at the intersection of wealth inequality and macroeconomics. In my job market paper, I build and estimate a structural heterogeneous agent model with returns that differ across households. The core contribution is to show how heterogeneity in returns and income risk jointly shape the wealth distribution and consumption behavior. Relative to closely related work, my framework performs better at matching the lower and middle portions of the wealth distribution. This difference matters quantitatively. Models that focus primarily on returns heterogeneity can fit the top of the wealth distribution closely but tend to generate too much wealth among the bottom half of households. In my framework, income uncertainty generates a precautionary saving motive that produces genuinely low wealth households. This leads to a realistic distribution of marginal propensities to consume and aggregate consumption responses that align with empirical estimates.

A key feature of my research is the application of quantitative methods to real world policy and infrastructure questions. I work with large datasets, develop analytical tools, and produce written reports and presentations that translate complex results into accessible insights. I am comfortable explaining technical concepts to non technical audiences and tailoring analysis to decision relevant questions. This work requires precision, judgment, and strong communication, skills that map directly to applied energy analytics.

I am particularly interested in Demand Side Analytics because of its role in addressing pressing questions about the future of the electric grid and natural gas distribution in North America. The opportunity to work closely with senior staff to deliver high quality analyses and tools in a policy relevant setting is especially appealing. I value environments where analytical rigor directly informs real world decisions.

My economics training has provided a strong foundation in quantitative analysis, statistical reasoning, and economic modeling. I am comfortable working independently in a remote setting while collaborating effectively with team members. I take pride in producing clear, well documented work products and in meeting deadlines consistently.

I would welcome the opportunity to contribute my skills and perspective to Demand Side Analytics. Thank you for your time and consideration.

\closing{Sincerely,}

\end{letter}
\end{document}
